\documentclass[12pt]{article}
\usepackage{amsmath, amssymb, geometry, enumitem}
\geometry{margin=1in}
\setlength{\parskip}{0.6em}
\setlength{\parindent}{0pt}

\begin{document}

\begin{center}
\Large \textbf{ECON 300 -- Labour Economics} \vspace{.5cm}\\
\large \textbf{Solutions -- Quiz 4} \vspace{.5cm}\\
\large \textbf{March 31, 2026}
\end{center}


\section*{Part A: Multiple Choice Solutions}

\begin{enumerate}[label=\textbf{\arabic*.}]

\item \textbf{Correct answer: C}

\textbf{Question:} In the Oaxaca decomposition, the portion of the wage gap that is attributed to differences in observed characteristics such as schooling and experience is called what?

\textbf{Explanation:}

The Oaxaca decomposition splits the total wage gap into two parts:

\begin{itemize}
    \item the \textbf{explained component}, which comes from differences in observable characteristics such as education, experience, and occupation;
    \item the \textbf{unexplained component}, which comes from differences in coefficients or returns to those characteristics.
\end{itemize}

Thus, the portion due to measured characteristics is called the \textbf{explained component}.

\[
\boxed{\text{Answer: C}}
\]

\item \textbf{Correct answer: C}

\textbf{Question:} Suppose women and men have the same observed productivity characteristics, but women still receive lower wages because firms statistically expect greater labour force interruptions among women. This is best described as what?

\textbf{Explanation:}

This is an example of \textbf{statistical discrimination}. Employers are not judging an individual woman solely on her own expected productivity. Instead, they use group averages or stereotypes about women's labour force attachment and apply them to all women.

This differs from:
\begin{itemize}
    \item \textbf{taste-based discrimination}, where employers simply dislike hiring women;
    \item \textbf{human capital investment}, which concerns education and training decisions;
    \item \textbf{comparable worth}, which is a policy concept.
\end{itemize}

\[
\boxed{\text{Answer: C}}
\]

\item \textbf{Correct answer: B}

\textbf{Question:} According to the crowding hypothesis, if women are restricted to a narrow set of occupations, what happens?

\textbf{Explanation:}

The \textbf{crowding hypothesis} says that if one group is excluded from many occupations, it becomes concentrated in fewer occupations. This raises labour supply in those occupations, which pushes wages downward.

So if women are crowded into a narrow set of jobs:

\begin{itemize}
    \item labour supply in those jobs rises;
    \item equilibrium wages in those jobs fall.
\end{itemize}

That is the core mechanism of the crowding argument.

\[
\boxed{\text{Answer: B}}
\]

\item \textbf{Correct answer: A}

\textbf{Question:} Which of the following would most likely reduce the unexplained portion of the gender wage gap in an Oaxaca decomposition?

\textbf{Explanation:}

The unexplained component becomes smaller when the regression includes more relevant variables that explain wage differences.

If schooling differs across groups, then \textbf{adding education} as a regressor moves part of the wage gap from the unexplained component to the explained component.

The other options do not help:
\begin{itemize}
    \item removing controls makes omitted variable bias worse;
    \item increasing sample size improves precision, but does not necessarily reduce unexplained differences;
    \item omitting experience worsens the model;
    \item using nominal instead of real wages does not address decomposition.
\end{itemize}

\[
\boxed{\text{Answer: A}}
\]

\item \textbf{Correct answer: D}

\textbf{Question:} Which of the following best describes the \textbf{unexplained} component in an Oaxaca decomposition?

\textbf{Explanation:}

The unexplained component is the part of the wage gap that remains after controlling for differences in observed characteristics. It arises because men and women receive different \textbf{returns} to the same characteristics.

For example, if an extra year of schooling raises male wages more than female wages, that contributes to the unexplained component.

Thus, the unexplained portion reflects \textbf{differences in coefficients or returns to characteristics}.

\[
\boxed{\text{Answer: D}}
\]

\end{enumerate}

\vspace{1em}

\section*{Part B: Long Numerical Problem Solutions}

We are given the following estimated earnings equations:

\[
\ln W_m = 1.20 + 0.09S + 0.05EXP
\]

\[
\ln W_f = 0.95 + 0.07S + 0.04EXP
\]

where:

\begin{itemize}
    \item \(W_m\) = annual earnings of men
    \item \(W_f\) = annual earnings of women
    \item \(S\) = years of schooling
    \item \(EXP\) = years of labour market experience
\end{itemize}

Average characteristics:

\[
\bar{S}_m = 16, \qquad \bar{EXP}_m = 12
\]

\[
\bar{S}_f = 15, \qquad \bar{EXP}_f = 10
\]

\subsection*{(a) Predicted log earnings for men}

Substitute male averages into the male wage equation:

\[
\ln W_m = 1.20 + 0.09(16) + 0.05(12)
\]

Compute each part:

\[
0.09(16) = 1.44
\]

\[
0.05(12) = 0.60
\]

So:

\[
\ln W_m = 1.20 + 1.44 + 0.60
\]

\[
\ln W_m = 3.24
\]

\[
\boxed{\ln W_m = 3.24}
\]

\subsection*{Interpretation}

This means the predicted \textbf{log annual earnings} for the average man in the sample are 3.24.

\vspace{0.8em}

\subsection*{(b) Predicted log earnings for women}

Substitute female averages into the female wage equation:

\[
\ln W_f = 0.95 + 0.07(15) + 0.04(10)
\]

Compute each part:

\[
0.07(15) = 1.05
\]

\[
0.04(10) = 0.40
\]

So:

\[
\ln W_f = 0.95 + 1.05 + 0.40
\]

\[
\ln W_f = 2.40
\]

\[
\boxed{\ln W_f = 2.40}
\]

\subsection*{Interpretation}

This means the predicted \textbf{log annual earnings} for the average woman in the sample are 2.40.

\vspace{0.8em}

\subsection*{(c) Raw gender gap in predicted log earnings}

The raw gender gap is:

\[
\ln W_m - \ln W_f
\]

Substitute the answers from parts (a) and (b):

\[
3.24 - 2.40 = 0.84
\]

\[
\boxed{\ln W_m - \ln W_f = 0.84}
\]

\subsection*{Interpretation}

The predicted male-female gap in \textbf{log earnings} is 0.84. Since these are log wages, this is approximately a large percentage gap in earnings.

\vspace{0.8em}

\subsection*{(d) Explained portion of the gap using the male wage structure}

Using the male coefficients as the reference group:

\[
(\bar{X}_m - \bar{X}_f)\hat{\beta}_m
\]

We compare differences in characteristics:

\[
\bar{S}_m - \bar{S}_f = 16 - 15 = 1
\]

\[
\bar{EXP}_m - \bar{EXP}_f = 12 - 10 = 2
\]

Using the male returns:

\[
\text{Explained gap} = 0.09(1) + 0.05(2)
\]

\[
= 0.09 + 0.10
\]

\[
= 0.19
\]

\[
\boxed{\text{Explained portion} = 0.19}
\]

\subsection*{Interpretation}

Out of the total gap of 0.84, an amount of 0.19 is explained by differences in observed characteristics:

\begin{itemize}
    \item men have one more year of schooling on average;
    \item men have two more years of experience on average.
\end{itemize}

\vspace{0.8em}

\subsection*{(e) Unexplained portion of the gap}

The unexplained portion is:

\[
\text{Unexplained gap} = \text{Total gap} - \text{Explained gap}
\]

\[
= 0.84 - 0.19
\]

\[
= 0.65
\]

\[
\boxed{\text{Unexplained portion} = 0.65}
\]

\subsection*{Interpretation}

This means most of the observed gap remains unexplained after controlling for schooling and experience.

\vspace{0.8em}

\subsection*{(f) Economic interpretation of the explained and unexplained components}

\textbf{Explained component:}

The explained part of the wage gap arises because men and women differ in their measurable characteristics. In this problem:

\begin{itemize}
    \item men have slightly more schooling;
    \item men have more labour market experience.
\end{itemize}

These differences account for part of the wage gap.

\textbf{Unexplained component:}

The unexplained part reflects differences in the \textbf{returns to characteristics} or other omitted factors. It may arise from:

\begin{itemize}
    \item labour market discrimination,
    \item unobserved ability differences,
    \item omitted occupational characteristics,
    \item different pay structures for men and women.
\end{itemize}

Thus:

\[
\boxed{\text{Explained gap} = 0.19, \qquad \text{Unexplained gap} = 0.65}
\]

The unexplained portion is much larger, suggesting that differences in schooling and experience explain only a small part of the total predicted earnings gap.

\vspace{1em}
\newpage
\section*{Final Answers Summary}

\subsection*{Multiple Choice}
\begin{center}
\begin{tabular}{ccccc}
1. C & 2. C & 3. B & 4. A & 5. D
\end{tabular}
\end{center}

\subsection*{Numerical Problem}
\begin{align*}
(a)\quad & \ln W_m = 3.24 \\
(b)\quad & \ln W_f = 2.40 \\
(c)\quad & \ln W_m - \ln W_f = 0.84 \\
(d)\quad & \text{Explained component} = 0.19 \\
(e)\quad & \text{Unexplained component} = 0.65
\end{align*}

\end{document}